% ભાગ: first-order-logic
% પ્રકરણ: models-theories
% વિભાગ: theories

\documentclass[../../../include/open-logic-section]{subfiles}

\begin{document}

\olfileid{fol}{mat}{the}

\olsection{પ્રથમ-ક્રમ સિદ્ધાંતોનાં ઉદાહરણો}

\begin{ex}
ભાષા~$\Lang L_<$માં કડક રેખીય ક્રમોનો સિદ્ધાંત નીચેના ગણ વડે
સ્વયંસિદ્ધીકૃત થાય છે:
\begin{align*}
\{\quad & \lforall[x][\lnot x < x], \\
& \lforall[x][\lforall[y][((x < y \lor y <
    x) \lor x = y)]], \\
& \lforall[x][\lforall[y][\lforall[z][((x < y
      \land y < z) \lif x < z)]]] \quad \}
\end{align*}
તે અભિપ્રેત !!{structure}sને સંપૂર્ણપણે વર્ણવે છે: દરેક કડક રેખીય
ક્રમ આ સ્વયંસિદ્ધિ-તંત્રનું નિદર્શ છે; અને ઊલટું, જો~$R$ ગણ~$X$ પર
રેખીય ક્રમ હોય, તો $\Domain M=X$ અને $\Assign{<}{M}=R$ ધરાવતી
સંરચના~$\Struct M$ આ સિદ્ધાંતનું નિદર્શ છે.
\end{ex}

\begin{ex}
$\Obj 1$ (!!{constant}) અને $\cdot$ (દ્વિસ્થાની !!{function}) ધરાવતી
ભાષામાં સમૂહોનો સિદ્ધાંત નીચે પ્રમાણે સ્વયંસિદ્ધીકૃત થાય છે:
\begin{align*}
& \lforall[x][\eq[(x \cdot \Obj 1)][x]]\\
& \lforall[x][\lforall[y][\lforall[z][\eq[(x \cdot (y \cdot z))][((x
          \cdot y) \cdot z)]]]]\\
& \lforall[x][\lexists[y][\eq[(x \cdot y)][\Obj 1]]]
\end{align*}
\end{ex}

\begin{ex}
અંકગણિતની ભાષા~$\Lang L_A$માં પેયાનો અંકગણિતનો સિદ્ધાંત નીચેના
વાક્યો વડે સ્વયંસિદ્ધીકૃત થાય છે:
\begin{align*}
& \lforall[x][\lforall[y][(\eq[x'][y'] \lif \eq[x][y])]]\\
& \lforall[x][\eq/[\Obj 0][x']]\\
& \lforall[x][\eq[(x + \Obj 0)][x]]\\
& \lforall[x][\lforall[y][\eq[(x + y')][(x + y)']]]\\
& \lforall[x][\eq[(x \times \Obj 0)][\Obj 0]]\\
& \lforall[x][\lforall[y][\eq[(x \times y')][((x \times y) + x)]]]\\
& \lforall[x][\lforall[y][(x < y \liff \lexists[z][\eq[(z' + x)][y])]]]\\
\intertext{અને આ સ્વરૂપનાં બધાં વાક્યો:}
& (!A(\Obj 0) \land \lforall[x][(!A(x) \lif !A(x'))]) \lif \lforall[x][!A(x)]
\end{align*}
છેલ્લા સ્વરૂપનાં વાક્યો અનંત સંખ્યામાં હોવાથી આ સ્વયંસિદ્ધિ-તંત્ર
અનંત છે. છેલ્લા સ્વરૂપને \emph{અનુમાનપ્રવર્તન પ્રરૂપ} કહે છે. (હકીકતમાં,
અનુમાનપ્રવર્તન પ્રરૂપને આપણે અહીં બતાવ્યું તેના કરતાં થોડું વધુ જટિલ છે.)

છેલ્લું સ્વયંસિદ્ધિ~$<$ની \emph{સ્પષ્ટ વ્યાખ્યા} છે.
\end{ex}

\begin{ex}
શુદ્ધ ગણોનો સિદ્ધાંત ગણિતના પાયામાં (અને ગણિતના તત્ત્વજ્ઞાનમાં) અગત્યની
ભૂમિકા ભજવે છે. કોઈ ગણના બધા !!{element}s પણ શુદ્ધ ગણો હોય તો તે ગણ
શુદ્ધ છે. તેથી ખાલી ગણ શુદ્ધ ગણાય છે, પરંતુ કોઈ ગણમાં પોતે ગણ ન હોય
એવી કોઈ વસ્તુ !!a{element} તરીકે આવે તો તે શુદ્ધ નથી. એટલે શુદ્ધ ગણો એવા ગણો છે જે માત્ર ખાલી ગણમાંથી
બને છે અને જેમાં કોઈ ``અગણ મૂળઘટક'' નથી, એટલે કે જે વસ્તુઓ પોતે ગણ
નથી એવી કોઈ વસ્તુ નથી.

નીચેનાને શુદ્ધ ગણોના સિદ્ધાંત માટેનું સ્વયંસિદ્ધિ-તંત્ર ગણી શકાય:
\begin{align*}
& \lexists[x][\lnot \lexists[y][y \in x]]\\
& \lforall[x][\lforall[y][(\lforall[z][(z \in x \liff z \in y)] \lif
\eq[x][y])]]\\
& \lforall[x][\lforall[y][\lexists[z][\lforall[u][(u \in z \liff
(\eq[u][x] \lor \eq[u][y]))]]]]\\
& \lforall[x][\lexists[y][\lforall[z][(z \in y \liff \lexists[u][(z \in
        u \land u \in x)])]]]\\
\intertext{અને આ સ્વરૂપનાં બધાં વાક્યો:} &
\lexists[x][\lforall[y][(y \in x \liff !A(y))]]
\end{align*}
\sourcecorrection{OLMAT-001}{મૂળની \texttt{theories.tex},
પંક્તિ~75માં ઘટકો દ્વારા નિર્ધારિત સમાનતાના સ્વયંસિદ્ધિમાં
\texttt{\string\lforall[z]} પછી દલીલનું આરંભક ચોરસ કૌંસ નથી અને તેના
સ્થાને ગોળ કૌંસ આવે છે. અહીં બીજા સ્વયંસિદ્ધિને સુમેળવાળા
$\lforall[z][\dots]$ સ્વરૂપે પુનઃસ્થાપિત કર્યું છે.}
પહેલું સ્વયંસિદ્ધિ કહે છે કે કોઈ !!{element}s ન ધરાવતો ગણ છે (એટલે
$\emptyset$ અસ્તિત્વ ધરાવે છે); બીજું કહે છે કે ગણો ઘટકો દ્વારા
નિર્ધારિત છે; ત્રીજું કહે છે કે કોઈ પણ ગણો~$X$ અને~$Y$ માટે ગણ
$\{X,Y\}$ અસ્તિત્વ ધરાવે છે; અને ચોથું કહે છે કે કોઈ પણ ગણ~$X$ માટે
ગણ~$\cup X$ અસ્તિત્વ ધરાવે છે, જ્યાં~$\cup X$ એટલે~$X$ના બધા ઘટકોનો
યોગગણ.

છેલ્લે જણાવેલાં !!{sentence}sને સામૂહિક રીતે \emph{સહજ ગુણધર્મ-ગણરચના
પ્રરૂપ} કહે છે. તેનો સાર એ છે કે દરેક~$!A(x)$ માટે ગણ
$\Setabs{x}{!A(x)}$ અસ્તિત્વ ધરાવે છે---એટલે પ્રથમ નજરે તે સાચું,
ઉપયોગી અને કદાચ જરૂરી પણ લાગતું સ્વયંસિદ્ધિ છે. તેને ``સહજ'' કહે છે,
કારણ કે હકીકતમાં તે આ સિદ્ધાંતને અસંતોષ્ય બનાવે છે: જો~$!A(y)$ને
$\lnot y\in y$ લો, તો આપણને !!{sentence}
\[
\lexists[x][\lforall[y][(y \in x \liff \lnot y \in y)]]
\]
મળે છે, અને આ !!{sentence} કોઈ પણ !!{structure}માં સંતોષાતું નથી.
\end{ex}

\begin{ex}
\emph{અંશસિદ્ધાંત} ક્ષેત્રમાં \emph{અંશતા}નો સંબંધ પાયાનો
સંબંધ છે. ગણોના સિદ્ધાંતોની જેમ અંશતાના પણ વિવિધ---અને કેટલીક વાર
પરસ્પર વિરોધી---અર્થઘટનોને સ્વયંસિદ્ધીકૃત કરતા સિદ્ધાંતો છે.

અંશસિદ્ધાંતની ભાષામાં માત્ર એક દ્વિસ્થાની વિધેય-સંકેત~$\Obj P$ છે,
અને $\Atom{\Obj P}{x,y}$નો ``અર્થ'' એ છે કે~$x$,~$y$નો અંશ છે. આ
અર્થઘટન મનમાં રાખીએ ત્યારે આ ભાષા માટેની !!a{structure}ને
\emph{અંશતા-સંરચના} કહે છે. અલબત્ત, માત્ર એક દ્વિસ્થાની વિધેય ધરાવતી
દરેક સંરચના ખરેખર આ નામને લાયક નથી. ``અંશતા''ને વર્ણવવાની શક્યતા માટે
$\Assign{\Obj P}{M}$એ કેટલીક શરતો સંતોષવી જોઈએ; અંશતાના સિદ્ધાંત માટે
તેમને સ્વયંસિદ્ધિઓ તરીકે મૂકી શકીએ છીએ. દાખલા તરીકે, અંશતા વસ્તુઓ પરનો
આંશિક ક્રમ છે: દરેક વસ્તુ પોતાનો અંશ (ભલે તે \emph{અયથાર્થ} અંશ હોય)
છે; બે જુદી વસ્તુઓ એકબીજાની અંશ હોઈ શકે નહિ; અને કોઈ વસ્તુના અંશનો
અંશ એ વસ્તુનો અંશ હોય છે. નોંધો કે આ અર્થમાં ``નો અંશ છે'' એ ``નો
ઉપગણ છે''ને મળતું આવે છે, પરંતુ સ્વવાચક કે પરંપરિત ન હોય એવા ``નો ઘટક
છે''ને મળતું આવતું નથી.
\begin{align*}
& \lforall[x][\Atom{\Obj P}{x,x}] \\
& \lforall[x][\lforall[y][((\Part{x}{y} \land \Part{y}{x})
      \lif \eq[x][y])]] \\
& \lforall[x][\lforall[y][\lforall[z][((\Part{x}{y} \land
        \Part{y}{z}) \lif \Part{x}{z})]]]\\
\intertext{વળી, કોઈ પણ બે વસ્તુઓનો અંશસિદ્ધાંતિક યોગ હોય છે (એવી વસ્તુ
  જેમાં આ બંને વસ્તુઓ અંશ તરીકે હોય અને આ બાબતમાં જે લઘુતમ હોય).} &
\lforall[x][\lforall[y][\lexists[z][\lforall[u][(\Part{z}{u} \liff
        (\Part{x}{u} \land \Part{y}{u}))]]]]
\end{align*}
આ તો તત્ત્વચિંતકોએ વિચારેલા અંશતાના પાયાના સિદ્ધાંતોમાંના માત્ર કેટલાક
છે. જોકે પહેલાં કેટલાક વ્યાખ્યાયિત સંબંધો દાખલ કર્યા વિના આગળના
સિદ્ધાંતોને ઘડવા કે લખવા ઝડપથી મુશ્કેલ બને છે. દાખલા તરીકે,
અંશસિદ્ધાંતમાં રસ ધરાવતા મોટા ભાગના તત્ત્વચિંતકો આને પણ પ્રમાણભૂત
સિદ્ધાંત ગણે છે: જ્યારે કોઈ વસ્તુ~$x$નો યથાર્થ અંશ~$y$ હોય, ત્યારે
તેનો એવો અંશ~$z$ પણ હોય કે તેનો અને~$y$નો કોઈ સામાન્ય અંશ ન હોય અને
$y$ તથા~$z$નું સંલયન~$x$ હોય.
\end{ex}

\end{document}
